\section{Recovery of Visible and Completed Geometry}
\label{sec:statistical-recovery}

The isometric fixed-channel section transfers visible SPD error exactly.  When
the information subspace is estimated, the stratified line element adds a
second, mismatch-weighted Grassmann term.

All results below are stated after congruence whitening the reference to
\(P_0=I\); AIRM congruence invariance transports them to an arbitrary fixed
reference.  In these coordinates, for an orthonormal frame
\(U\in\St(d,m)\), write
\begin{equation}
  \Cpl(U,R)=I+U(R-I)U^\top.
  \label{eq:normalized-completion}
\end{equation}

Two observation regimes must be distinguished.  The exact visible space in
\cref{sec:introduction} is the span of the observed covectors.  The
known-subspace result below conditions on any such fixed space.  The
unknown-subspace result instead adopts a population second-moment model and
\emph{defines} its target space as the leading \(m\)-dimensional eigenspace
of \(\Sigma\).  This spectral target agrees with the population covector
support under a rank-\(m\) population model; a finite observed span equals it
only when the sampled covectors span that support.  With nonzero spectral
tails it instead represents a chosen principal approximation.  No automatic
identification of the two regimes is assumed.

For \(W,\widetilde W\in\Gr(m,d)\), let \(\dgr(W,\widetilde W)\) be the
Euclidean norm of their principal-angle vector.  Define
\begin{equation}
  \chi(R)=\opnorm{\Xiop(R)}^{1/2}
  =\max_{\lambda\in\operatorname{spec}(R)}
  \frac{|\lambda-1|}{\sqrt\lambda}.
  \label{eq:chi}
\end{equation}

\begin{theorem}[Joint deterministic perturbation]
\label{thm:joint-perturbation}
Let \(U,V\) be endpoint frames connected by a horizontal minimizing
Grassmann geodesic, and express \(R,S\in\SPD^m\) in the parallel endpoint
frames.  With reference \(I\),
\begin{equation}
  \boxed{
  \dai(\Cpl(U,R),\Cpl(V,S))
  \le
  \dai(R,S)
  +\sqrt2\min\{\chi(R),\chi(S)\}\dgr(W,\widetilde W).}
  \label{eq:joint-perturbation}
\end{equation}
The first term is visible geometry error.  The second is subspace error, and
it vanishes for reference-valued visible geometry.
\end{theorem}

The endpoint-frame condition is a gauge choice, not an additional statistical
assumption.  Given arbitrary orthonormal endpoint frames, align them by the
orthogonal Procrustes factor associated with the minimizing Grassmann geodesic
and conjugate the endpoint visible matrix accordingly.  The statement then
applies in the aligned coordinates, and its value is independent of the
initial endpoint gauges.

For a known fixed visible subspace, a direct concentration theorem retains
the intrinsic dimension.  In an orthonormal coordinate frame of \(W\), let
\(X_1,\ldots,X_n\in\R^m\) be i.i.d. copies of \(X\), define the uncentered
second moment \(R=\E[XX^\top]\succ0\), assume
\(\norm{R^{-1/2}X}_2\le L\) almost surely, and set
\(\Rhat=n^{-1}\sum_iX_iX_i^\top\).

\begin{theorem}[Known-subspace finite-sample recovery]
\label{thm:known-subspace-recovery}
For \(0<\varepsilon\le1\), if
\begin{equation}
  n\ge
  \frac{8L^4}{\varepsilon^2}
  \left(m\log9+\log\frac{2}{\delta}\right),
  \label{eq:known-subspace-sample-size}
\end{equation}
then with probability at least \(1-\delta\),
\begin{equation}
  e^{-\varepsilon}R\preceq\Rhat\preceq e^\varepsilon R,
  \qquad
  \dai(\Cpl(U,R),\Cpl(U,\Rhat))\le\sqrt m\,\varepsilon.
  \label{eq:known-subspace-bound}
\end{equation}
\end{theorem}

This self-contained net--Hoeffding bound is deliberately conservative.  In
fact
\[
  m=\E\norm{R^{-1/2}X}_2^2\le L^2,
\]
so the displayed sufficient sample size can have order
\(m^3/\varepsilon^2\) when only this envelope is retained.  Matrix
Chernoff or Bernstein inequalities can replace the net argument and often
improve the dependence to a logarithmic dimension or an effective-rank
factor under their corresponding variance or tail assumptions
\citep{tropp2015matrix}.  We use the elementary bound for transparency and
do not claim an optimal statistical rate.

We next allow the visible eigenspace to be unknown.  Let \(1\le m<d\) and
\(\Sigma\in\PSD^d\) have eigenvalues
\(\lambda_1\ge\cdots\ge\lambda_d\), assume
\(\gamma=\lambda_m-\lambda_{m+1}>0\), and let \(U\) span its leading
\(m\)-dimensional eigenspace.  Put
\[
  R=U^\top\Sigma U,
  \qquad
  P=\Cpl(U,R).
\]
For an estimator \(\Sigmahat\), let \(V\) span its leading eigenspace,
\(\Rhat=V^\top\Sigmahat V\), and \(\Phat=\Cpl(V,\Rhat)\).

\begin{theorem}[Unknown-subspace completed-geometry recovery]
\label{thm:unknown-subspace-recovery}
Let
\[
  \opnorm{\Sigmahat-\Sigma}\le\delta
  <\min\{\gamma/2,\lambda_m\},
  \qquad
  s_\delta=\frac{\delta}{\gamma-\delta},
\]
\[
  \varepsilon_{\mathrm{vis}}
  =\delta+2\sqrt2\lambda_1s_\delta,
  \qquad
  \eta_\delta=\frac{\varepsilon_{\mathrm{vis}}}{\lambda_m}<1,
\]
and
\[
  \omega_\delta
  =\max_{x\in[\lambda_m-\delta,\lambda_1+\delta]}
  \frac{|x-1|}{\sqrt x}.
\]
Then
\begin{equation}
  \boxed{
  \dai(P,\Phat)
  \le
  \sqrt m\left[-\log(1-\eta_\delta)
  +\sqrt2\,\omega_\delta\arcsin(s_\delta)\right].}
  \label{eq:unknown-subspace-bound}
\end{equation}
The first term controls the aligned visible second moment; the second is the
mismatch-weighted rotation of the estimated eigenspace.
\end{theorem}

\begin{corollary}[A self-contained sample bound]
\label{cor:unknown-subspace-sample}
Let \(X_1,\ldots,X_n\in\R^d\) be i.i.d. copies of \(X\), with uncentered
second moment \(\Sigma=\E[XX^\top]\), \(\norm X_2\le L\) almost surely,
and
\(\Sigmahat=n^{-1}\sum_iX_iX_i^\top\).  Define
\begin{equation}
  \delta_n(\alpha)
  =L^2\sqrt{\frac{2(d\log9+\log(2/\alpha))}{n}}.
  \label{eq:ambient-delta}
\end{equation}
With probability at least \(1-\alpha\),
\(\opnorm{\Sigmahat-\Sigma}\le\delta_n(\alpha)\).  If
\(\delta_n(\alpha)\) satisfies the hypotheses of
\cref{thm:unknown-subspace-recovery}, then
  \eqref{eq:unknown-subspace-bound} holds with
  \(\delta=\delta_n(\alpha)\).

In particular, put
\begin{equation}
  c_\gamma=1+\frac{8\sqrt2\lambda_1}{3\gamma},
  \qquad
  \delta_\star
  =\min\left\{\frac{\gamma}{4},
  \frac{\lambda_m}{2c_\gamma}\right\}.
  \label{eq:explicit-delta-threshold}
\end{equation}
The directly checkable sufficient condition
\begin{equation}
  n\ge
  \frac{2L^4}{\delta_\star^2}
  \left(d\log9+\log\frac{2}{\alpha}\right)
  \label{eq:explicit-unknown-sample-size}
\end{equation}
implies \(\delta_n(\alpha)\le\delta_\star\), hence
\(s_\delta\le4\delta/(3\gamma)\) and
\(\eta_\delta\le1/2\).  It therefore guarantees all hypotheses of
\cref{thm:unknown-subspace-recovery}.
\end{corollary}

The ambient dimension in \eqref{eq:ambient-delta} is the price of estimating
the subspace with a self-contained bounded-i.i.d. net argument.  Matrix
concentration or effective-rank assumptions may improve this factor, but an
eigengap or another identifiable subspace condition remains necessary.  The
reference \(I\) here is the whitened reference geometry, and \(\Sigma\) is an
uncentered second moment; it coincides with covariance only when the mean is
zero.  Proofs are in
\cref{app:recovery-proofs}.
